\documentclass[12pt,a4paper]{report}
\usepackage[utf8]{inputenc}
\usepackage{amsmath}
\usepackage{amsfonts}
\usepackage{amssymb}
\usepackage{amsthm}
\usepackage{hyperref}

\usepackage{multicol}
\usepackage{fancyhdr}
\usepackage[inline]{enumitem}
\usepackage{tikz}
\usepackage{tikz-cd}
\usetikzlibrary{calc}
\usetikzlibrary{shapes.geometric}
\usetikzlibrary{positioning}
\usepackage[margin=0.5in]{geometry}
\usepackage{xcolor}

\hypersetup{
    colorlinks=true,
    linkcolor=blue,
    filecolor=magenta,      
    urlcolor=cyan,
    pdftitle={Tensors},
    pdfpagemode=FullScreen,
    }

%\urlstyle{same}

\newcommand{\CLASSNAME}{Math 5050 -- Special Topics: Differential Geometry}
\newcommand{\STUDENTNAME}{Paul Carmody}
\newcommand{\ASSIGNMENT}{Section 21: Orientations }
\newcommand{\DUEDATE}{October 15, 2025}
\newcommand{\PROFESSOR}{Professor Berchenko-Kogan}
\newcommand{\SEMESTER}{Fall 2025}
\newcommand{\SCHEDULE}{TBD}
\newcommand{\ROOM}{Remote}

\newcommand{\MMN}{M_{m\times n}}
\newcommand{\FF}{\mathcal{F}}

\pagestyle{fancy}
\fancyhf{}
\chead{ \fancyplain{}{\CLASSNAME} }
%\chead{ \fancyplain{}{\STUDENTNAME} }
\rhead{\thepage}
\input{../MyMacros}

\newcommand{\RED}[1]{\textcolor{red}{#1}}
\newcommand{\BLUE}[1]{\textcolor{blue}{#1}}

\newcommand{\SUPP}{\operatorname{supp}}
\newcommand{\CLOSURE}{\operatorname{closure of}}

\begin{document}

\begin{center}
	\Large{\CLASSNAME -- \SEMESTER} \\
	\large{ w/\PROFESSOR}
\end{center}
\begin{center}
	\STUDENTNAME \\
	\ASSIGNMENT -- \DUEDATE\\
\end{center} 


\begin{enumerate}[label=21.\arabic*]

\item \textbf{Locally constant map on a connected space}

A map $f: S \to Y$ between two topological spaces is \textit{locally constant} if for every $p\in S$ there is a neighborhood $U$ of $p$  such that $f$ is constant on $U$.  Show that a locally constant map $f:S\to Y$ on a nonempty connected space $S$ is constant.  (\textit{Hint: }Show that every $y \in Y$, the inverse image $f^{-1}(y)$ is open.  Then $S =\cup_{y\in Y}f^{-1}(y)$ exhibits $S$ as a disjoint union of open subsets.)

\item \textbf{Continuity of pointwise orientations}

Prove that a pointwise orientation $[(X_1,\dots, X_n)]$ on a manifold $M$ is continuous if and only if every point $p \in M$  has a coordinate neighborhood $(U,\phi)=(U, x^1,\dots, x^n)$ such that for all $q \in U$, the differential $\phi_{*,q}:T_qM\to T_{f(q)}\R^n\cong\R^n$ carries the orientation of $T_qM$ to the standard orientation of $\R^n$ in the following sense: $(\phi_*X_{1,q},\dots,X_{n,q})\sim(\PART{}{r^1},\dots,\PART{}{r^n})$.

\item \textbf{Equivalence of oriented atleses}

Show that the relation in Definition 2.11 is an equivalence relation.

\BLUE{An Equivalance Relation is symmetric, Reflexive and Transitive
\begin{itemize}
	\item \textbf{Symmetric: $a\sim b \implies b \sim a$}\\
	Let $J_\leftarrow$ be the Jacobian for $\phi_\alpha \circ \psi^{-1}_\beta$ and $J_\rightarrow$ be the Jacobian for $\psi_\beta \circ \phi_\alpha^{-1}$.  Clearly, $J_\leftarrow J_\rightarrow = \mathbb{I}$.  Thus both $J_\leftarrow$ and $J_\rightarrow$ must have the same sign.
	\item \textbf{Reflexive: $a\sim a$}
	$\phi_\alpha \circ \phi_\alpha = \mathbb{I}$ whose Jacobian is positive.		
	\item \textbf{Transitive $a \sim b, b\sim c \implies a \sim c$}.\\
	Let $J_a$ be the Jacobian for $\phi_\alpha \circ \psi^{-1}_\beta$ and $J_b$ be the Jabobian for $\psi_\beta \circ \tau^{-1}_\gamma$ for some chart $(W, \tau_\gamma)$.  Then, $J_c$ the Jacobian from $(U, \phi_\alpha)$ to $(W, \tau_\gamma)$ would be $J_c = J_a J_b$.  Thus, $J_c$ will be the same sign as $J_a$ and $J_b$ if they also have the same sign.
\end{itemize}
}

\item \textbf{Orientation-preserving diffeomorphisms}

Let $F: (N,[\omega_N])\to (M,[\omega_M])$ be an orientation-preserving diffeomorphism.  If $\{(V_\alpha, \psi_\alpha)\} = \{(V_\alpha, y_\alpha^1,\dots, y_\alpha^n)\}$ is an oriented atlas on $M$ that specifies the orientation of $M$, show that $\{(F^{-1}V_\alpha, FV^*\psi_\alpha)\}=\{(F^{-1}V_\alpha,F_\alpha^1,\dots,F_\alpha^n)\}$ is an oriented atlas on $N$ that specifies the orientation of $N$, where $F_\alpha^i=y_\alpha\circ F$.

\item \textbf{Orientation-preserving or orientation-reversing diffeomorphisms}

Let $U$ be an open set $(0,\infty)\times (0,2\pi)$ in the $(r,\theta)$-plane $\R^2$.  We define $F:U\subset \R^2 \to \R^2$ by $F(r,\theta)=-(r\cos\theta,r\sin\theta)$.  Decide whether $F$ is orientation-preserving or orientation-reversing as a diffeomorphism onto its image.

\BLUE{
\begin{align*}
	x(r,\theta) &= -r\cos\theta \\
	y(r,\theta) &= -r\sin\theta \\
	\PART{x}{r} &= -\cos\theta  &\quad \PART{x}{\theta} &= r\cos\theta \\
	\PART{y}{r} &= -\sin\theta & \quad \PART{y}{\theta} &= -r\sin\theta \\
	J &= \TWOXTWO{-\cos\theta}{r\sin\theta}{-\sin\theta}{-r\cos\theta} \\
	\det J &= (-\cos\theta)(-r\cos\theta)-(r\sin\theta)(-\sin\theta) \\ 
	&= r^2
\end{align*}Clearly, $\det J >0$ therefore orientation-preserving.
}

\item \textbf{Orientability of a regular level set in $\R^{n+1}$}

Suppoe $f(x^1,\dots,x^{n+1})$ is a $C^\infty$ function on $\R^{n+1}$ with 0 as a regular value.  Show that the zero set of $f$ is an orientable submanifold of $R^{n+1}$.  In particular, the unit $n$-sphere $S^n$ in $\R^{n+1}$ is orientable.

\item \textbf{Orientablity of a Lie group}

Show that every Lie group $G$ is orientable by constructing a nowhere-vanishing top form on $G$.

\item \textbf{Orientability of a parallelizable manifold.}

Show that the parallelizable manifold is orientable.  (In particular, this shows again that every Lie group is orientable.)

\item \textbf{Orientability of the total space of the tangent bundle}

Let $M$ be a smooth manifold and $\pi:TM\to M$ its tangent bundle.  Show that if $\{(U,\phi)\}$ is any atlas on $M$, then the atlas $\{(TU,\hat{\phi})\}$ on $TM$, with $\hat{\phi}$ defined in equation (12.1), is oriented.  This proves that the total space $TM$ of the tangent bundle is always orientable, regardless of whether $M$ is orientable.

\item \textbf{Oriented atlas on a circle}

In Example 5.16 we found an atlas $\mathfrak{U}=\{(U_i,\phi_i)\}_{i=1}^4$ on the unit circle $S^1$.  Is ${U}$ an oriented atlas?  If not, alter the coordinate function $\phi_i$ to make $\mathfrak{U}$ into an oriented atlas.

\BLUE{
Recall that 
	\begin{align*}
		\mathfrak{U} = \BRACKET{\begin{array}{l}
		(U_1 = \{ (x,y) \in S^1: y > 0\},\, \phi_1(x,y)=x), \\
		(U_2 = \{ (x,y) \in S^1: y < 0\},\, \phi_2(x,y)=x), \\
		(U_3 = \{ (x,y) \in S^1: x > 0\},\, \phi_3(x,y)=y), \\
		(U_4 = \{ (x,y) \in S^1: x < 0\},\, \phi_4(x,y)=y)
		\end{array}
	}
	\end{align*}
Need to show that the determinant of the transformation matrices (Jacobian) between each of the charts is positive.  Keeping in mind that the transformation only occurs on the intersection of the charts.  Taking $\phi_\alpha\circ\phi_\beta^{-1}$ for $\alpha, \beta \in \{1,2,3,4\}$ and its inverse we see 
\begin{enumerate}
	\item $(U_1, \phi_1) \cap (U_2,\phi_2)$ is empty.
	\item $(U_1, \phi_1) \cap (U_3,\phi_3)$ is the first quadrant.
	\begin{align*}
		\phi_1 \circ \phi_3^{-1}(y) &= \phi_1(\phi_3^{-1}(y)) = \phi(\sqrt{1-y^2},1) = \sqrt{1-y^2}\\
		\phi_3\circ\phi_1^{-1}(x) &= \phi_3(\phi_1^{-1}(x)) = \phi(1,\sqrt{1-x^2},1) = \sqrt{1-x^2}
	\end{align*}
	\item $(U_1, \phi_1) \cap (U_4,\phi_4)$ is the second quadrant.
	\begin{align*}
		\phi_1\circ\phi_4^{-1}(y) &= \phi_1(\phi_4^{-1}(y)) = \phi_1(-\sqrt{1-y^2}, y) = -\sqrt{1-y^2} \\
		\phi_4\circ\phi_1^{-1}(x) &= \phi_4(\phi_1^{-1}(x)) = \phi_4(x,\sqrt{1-x^2}) = \sqrt{1-x^2}
	\end{align*}
	\item $(U_2, \phi_2) \cap (U_3,\phi_3)$ is the fourth quadrant
	\begin{align*}
		\phi_2\circ\phi_3^{-1}(y) &= \phi_2(\phi_3^{-1}(y)) = \phi_2(\sqrt{1-y^2}, y) = \sqrt{1-y^2} \\
		\phi_3\circ\phi_2^{-1}(x) &= \phi_3(\phi_2^{-1}(x)) = \phi_3(x, -\sqrt{1-x^2}) = -\sqrt{1-x^2}
	\end{align*}
	\item $(U_2, \phi_2) \to (U_4,\phi_4)$ is the third quadrant.
	\begin{align*}
		\phi_2\circ\phi_4^{-1}(y) &= \phi_2(\phi_4^{-1}(y)) = \phi_2(-\sqrt{1-y^2}, y) = -\sqrt{1-y^2} \\
		\phi_4\circ\phi_2^{-1}(x) &= \phi_4(\phi_2^{-1}(x)) = \phi_4(x, -\sqrt{1-x^2}) = -\sqrt{1-x^2}
	\end{align*}
	\item $(U_3, \phi_3) \to (U_4,\phi_4)$ is empty
\end{enumerate}We can see that in (c) and (d) the signs are different for the inverse of each composition.  Therefore, it is NOT orientation preserving.  However, if we change $\phi_2(x,y)=-x$ and $\phi_4(x,y) = -y$ this makes ALL of the transformations across all charts to have the same sign.  Hence, orientation preserving.
}

\end{enumerate}
\end{document}
